% 来源: 19c9fb91-7632-8e2b-8000-00001e8fecc6_2.7_可靠传输.pdf
% 提取时间: 2026-02-28T22:51:33+08:00

2.7 可靠传输
分布式系统显然是复杂的事物。它们存在缺乏同步性、缺乏确定性和缺乏信任等问题。因此在分布式系
统中，消息的接收时间可能差异很大；消息可能丢失，服务器可能崩溃，而且当一条消息到达时，它甚
至可能包含病毒。用Lamport那句著名的话来说，分布式系统就是 “在其中，一台你甚至都不知道存在的
计算机发生故障，都可能导致你自己的计算机无法使用”。
在数据传输过程中，可能会出现多种类型的差错，主要包括以下几种：
 ● 随机错误：也称为独立错，是指数据位流中的错误是随机发生的，各个错误之间相互独立，互
   不关联。这种错误通常是由信道中的热噪声等随机因素引起的。热噪声是一种持续存在的、具
   有随机性的噪声，它会在数据传输过程中对信号产生干扰，导致个别数据位发生错误。例如，
   在无线通信中，由于信号在传输过程中受到大气噪声、电磁干扰等因素的影响，可能会出现随
   机错。
 ● 突发差错：突发错是指在一段时间内，数据位流中出现连续的多个错误。这种错误通常是由突
   发的干扰事件引起的，如脉冲噪声、信号衰落、传输线路的瞬间故障等。例如，在有线通信
   中，当线路受到雷击、电磁脉冲干扰时，可能会导致一段连续的数据出现错误；数据率越高，
   突发性差错影响越大。在无线通信中，当信号遇到障碍物或受到多径效应影响时，也可能会出
   现突发错。
实际的数据传输过程中，如下图，往往是随机错和突发错同时存在，形成混合错。例如，在无线通
信环境中，既有热噪声引起的随机错，也有因信号衰落、干扰等因素导致的突发错。这种混合错的
存在增加了数据传输差错检测和纠正的复杂性。




交织技术可以将较长的突发错误转换为随机分布的错误，再结合使用如前向纠错码(FEC)等技术来
消除这些随机错误。因此，通常会联合使用信道编码与交织技术，以显著提升通信系统的整体可靠
性。交织的具体实现机制之一：矩阵交织。假设有一组由4比特组成的消息分组，从每个分组中提
取第一个比特，并将这四个比特组合成一个新的四比特帧。对于剩余的比特2至4，执行相同的操
作。接着，按照第一比特组成的帧、第二比特组成的帧……依次发送。如果在传输过程中某一帧
（例如第二个比特组成的帧）丢失，那么在没有应用交织的情况下，整个对应的消息分组将无法恢
                                                 1
复；而采用交织技术后，则仅有各消息分组中的特定位置（此处为第2比特）受到影响。假如所使
用的纠错编码方案具备修复任一四比特代码中单一错误的能力，则通过该信道编码手段，仍有可能
完整恢复所有分组内的信息。这就是交织技术的基本原理。简而言之，交织是通过将一个编码块中
的b个比特重新分配到n个帧中去改变它们之间的相对位置关系，理论上n值越大，系统的传输性能
越佳，但同时也会增加传输延迟。




总的来说，差错控制技术的种类包括
 ● 前向纠错 (Forward Error Correction)
 ● 检错重发 (error detection and retransmission) - ARQ
 ● ARQ+FEC -> Hybrid automatic repeat-request (HARQ)
    ○ 纠错和检错结合的工作方式简称纠检结合，自动在纠错和检错之间转换
    ○ 当错码数量少时，系统按前向纠错方式工作，以节省重发时间，提高传输效率；
    ○ 当错码数量多时，系统按反馈重发方式纠错，以降低系统的总误码率。




                                                       2
